% RESEARCH_DRAFT_SHA256: 763e01a8ca9e689ac2dba962af13efbd0f0ef8e0ffaac6e931f334ab22f3937a
% Layout and source-reference wrappers expanded for the web renderer.
\documentclass[11pt]{article}
\usepackage[margin=25mm]{geometry}
\usepackage{fontspec,unicode-math}
\setmainfont{Latin Modern Roman}
\setmathfont{Latin Modern Math}
\usepackage{amsmath,amssymb,mathtools,xurl,hyperref}
\hypersetup{colorlinks=true,urlcolor=blue,linkcolor=blue}
\setlength{\parindent}{0pt}
\setlength{\parskip}{5pt}
\setlength{\emergencystretch}{3em}
\newcommand{\sref}[2]{\hyperlink{source:#1:#2}{[S#2]}}
\newcommand{\cataloguescope}[1]{\paragraph{Source question.}#1}
\newcommand{\floor}[1]{\left\lfloor#1\right\rfloor}
\begin{document}

\section*{Top-k perfect matching: extendable subsets and a failed threshold minimax swap}
\noindent catalogue entry; catalogue entry. Research attempt dated 11 October 2026.
\paragraph{Result in scope.}This manuscript contains a source-specific attempt and complete proofs of the stated auxiliary calculations. It makes no new solution or novelty claim for the entire target.
\section{Original problem and preserved definition}
The following definition and source list are reproduced from the repository. Any earlier inspection dates and reported literature claims in this preserved text describe that source record, not additional verification by this attempt.
% BEGIN PRESERVED DEFINITION
\subsection{Definitions and mathematical statement}

\paragraph{Shared notation.}Standard number systems are denoted by $\mathbb N,\mathbb Z,\mathbb Q,\mathbb R,\mathbb C$; index ranges, computational models and specialized notation follow the cited problem source and the definitions below.


A perfect matching is a set of disjoint edges covering every vertex exactly once. If its edge weights in nonincreasing order are $w_1\geq\cdots\geq w_t$, its top-$k$ objective is $\sum_{i=1}^kw_i$, with $0\leq k\leq t$. The input weights are finitely encoded, and polynomial time is measured in their binary length and the graph size.
\paragraph{Statement recorded in the supplied source.}
Given a weighted graph and an integer $k$, is there a deterministic polynomial-time algorithm that finds a perfect matching maximizing the total weight of its $k$ heaviest edges? \href{https://doi.org/10.4171/owr/2021/53}{[S2]}
\paragraph{Scope qualification.}
The question asks for a deterministic algorithm on general graphs, and does not restrict to bipartite graphs. If no perfect matching exists the algorithm must detect infeasibility.
\paragraph{Source status.}
The supplied catalogue labels this record partially solved. Its source-status metadata was checked against the supplied snapshot on 2026-10-11. This import does not independently establish current open status or mathematical correctness. The catalogue literature assessment is dated 2026-08-22.
\subsection{Short English statement}
Find a perfect matching whose $k$ heaviest edges have the largest possible total weight.
\subsection{Sources}
\begin{itemize}
\item[S1] Ulam AI and the UnsolvedMath Contributors, UnsolvedMath: A Curated Collection of Open Mathematics Problems, catalogue entry. Catalogue curation and metadata: CC BY 4.0. The supplied source question is retained; definitions and the English summary are editorial additions.
\url{https://huggingface.co/datasets/ulamai/UnsolvedMath}.
\item[S2] Combinatorial Optimization (2022), Oberwolfach Reports, DOI 10.4171/owr/2021/53, pp. 2893--2954
\url{https://doi.org/10.4171/owr/2021/53}.
\item[S3] Combinatorial Optimization, Oberwolfach Reports 53/2021, primary report PDF
\url{https://ems.press/content/serial-article-files/46931}.
\textit{Underlying Oberwolfach Reports material retains its source-specific licence, including CC BY-SA 4.0 where applicable; see THIRD\_PARTY\_NOTICES.md.}
\end{itemize}
% END PRESERVED DEFINITION
\subsection{Verbatim supplied catalogue question}
Given a weighted graph and an integer $k$, is there a deterministic polynomial-time algorithm that finds a perfect matching maximizing the total weight of its $k$ heaviest edges?

\section{The source graph class and the catalogue target}
The primary passage explicitly starts with a bipartite graph, whereas
the catalogue asks about weighted graphs generally. The identities
and counterexample below apply to either class; the counterexample
itself is bipartite. The explicit matching-feasibility routine used
for the fixed-$k$ enumeration is justified for bipartite graphs.
No unrestricted deterministic algorithm for varying $k$ is claimed.

\paragraph{Exact equivalence with extendable partial matchings.}
Suppose a graph has a perfect matching with $t$ edges and
$0\le k\le t$. Let $\mathcal R_k$ be the $k$-edge
sets contained in some perfect matching. Then
\[
\max_M\operatorname{top}_k(M)=
\max_{R\in\mathcal R_k}\sum_{e\in R}w_e.
\]
For any perfect matching, its $k$ heaviest edges belong to
$\mathcal R_k$, proving one inequality. Conversely any $R$
in that family extends to a matching $M$, and the sum of
the largest $k$ weights in $M$ is at least the sum on
that particular $k$-subset, proving the other. Negative weights
and ties cause no change to these two comparisons. For $k=0$
both objectives are zero once feasibility is established.

For fixed $k$, enumerate all $k$-subsets of edges, discard
those sharing an endpoint, delete the endpoints of each remaining
set and test whether the residual graph has a perfect matching.
This is exactly the extension test: its union with $R$ is
perfect if and only if the residual matching is perfect. In a
bipartite graph a maximum matching is found by repeatedly augmenting
along alternating paths. Correctness follows because the symmetric
difference with any larger matching has an augmenting path component;
each search and the at most $|V|/2$ augmentations have polynomial
cost. Hence enumeration takes $|E|^{O(k)}$ time and is
polynomial for every fixed $k$. It is not a polynomial bound
when $k$ is part of the input, since the exponent depends on $k$.

\paragraph{The threshold identity and the tempting illegal exchange.}
For real numbers $a_1\ge\cdots\ge a_t$ and $1\le k<t$,
\[
\sum_{i=1}^k a_i
=\min_{\theta\in\mathbb R}
\left(k\theta+\sum_{i=1}^t\max\{a_i-\theta,0\}\right).
\]
For any $\theta$, each of the first $k$ terms is at most
$\theta+\max\{a_i-\theta,0\}$; summing and including
the remaining nonnegative terms gives the upper-bound direction.
If $a_{k+1}\le\theta\le a_k$, the first $k$ positive
parts sum to $\sum_{i\le k}(a_i-\theta)$ and all remaining
parts vanish, proving equality. This identity for one matching
does not allow interchanging the maximum over matchings with this
minimum over thresholds.

\paragraph{An exact bipartite minimax gap.}
Take a six-cycle. It has exactly two perfect matchings, its two
alternating triples: choosing the edge covering one vertex forces
successive choices around the cycle, giving precisely these two.
Give one triple weights $(10,0,0)$ and the other $(6,6,6)$,
and set $k=2$. Their top-two values are ten and twelve, so
the true optimum is twelve. For $0\le\theta\le6$ their
threshold objectives are $10+\theta$ and $18-\theta$.
The maximum is minimized at $\theta=4$ with value fourteen.
For $\theta<0$ the second objective is $18-\theta>18$;
for $6\le\theta\le10$ the first is $10+\theta\ge16$;
and for $\theta\ge10$ both are $2\theta\ge20$.
Consequently
\[
\max_M\min_\theta\Phi(M,\theta)=12
\quad\hbox{but}\quad
\min_\theta\max_M\Phi(M,\theta)=14.
\]
The inner maximization at fixed threshold is an ordinary weighted
matching objective, but solving that family of objectives optimizes
the displayed upper bound, not an equal reformulation of the
original problem. This explicit gap blocks that proposed proof of
a deterministic solution; it does not exclude other threshold-based
algorithms with an additional argument.

\section{Audit}
The original top-$k$ passage on printed page 2948 was inspected.
The audit corrected the source/catalogue graph-class discrepancy,
retained binary-encoded unrestricted weights, and did not equate
fixed-$k$ polynomial time with a uniform polynomial algorithm.
The six-cycle calculation is exact hand enumeration of two matchings;
no matching solver or random isolation procedure was run.
Derandomization or another uniform deterministic optimization mechanism
for varying $k$ remains missing in this attempt.

\section{Process and contributions}
Mathematics and model-assisted adversarial self-audit: OpenAI gpt-6-astra, ultra reasoning, as dispatched for this research batch; exact serving revision not exposed. Codex (GPT-6; exact serving revision not exposed) materially assisted with orchestration and format review. No human mathematical authorship, mathematical edits, or independent human proof checking is asserted.
Contribution prompt A was applied to the full definition and supplied source row, then prompt B was applied adversarially to the derived results. The item-specific audit above records the scope checks and rejected extrapolations. Local assembly and structural checks do not establish mathematical correctness. The source catalogue status is preserved in TOP\_PROBLEM and is separate from this attempt status.
\paragraph{Integration review.} Codex (GPT-6; exact serving revision not exposed) reviewed the written arguments and expanded layout and source-reference wrappers. Repository metadata and the production web renderer were checked using Python 3.13 and JavaScript with MathJax 3.2.2. This model-assisted review is not independent human proof verification; no PDF compilation or formal proof checking is asserted.

\section{Completion Estimate}
\noindent\textbf{COMPLETION: 0\%}
The proved material is a rederivation of known or auxiliary facts rather than a substantiated new advance on the original target. Exact extendable-subset equivalence and fixed-k bipartite enumeration do not give uniform polynomial time; a concrete minimax gap defeats the naive threshold interchange.
\section{References}
The complete original Sources subsection above is retained, including catalogue attribution and licence notices. The additional list records primary sources actually consulted for this attempt on 11 October 2026; uncited source links reproduced above were not all rechecked.
\begin{itemize}
\item Oberwolfach Report 53/2021, Combinatorial Optimization, original question and surrounding discussion read from the primary PDF text on 11 October 2026. Historical statements are not a current-status survey. \url{https://ems.press/content/serial-article-files/46931}.
\end{itemize}
\section{Final status}
\textbf{Unresolved in this attempt.} Exact extendable-subset equivalence and fixed-k bipartite enumeration do not give uniform polynomial time; a concrete minimax gap defeats the naive threshold interchange.
\end{document}
